\subsection{Large corners with four objects and fifteen arrows}

\begin{proposition}\label{C:fifteen-four-large-corner}
A four-object category with at most fifteen arrows, one of whose
endomorphism monoids has at least eight elements, cannot realize
$\mathbb{CP}^2$.
\end{proposition}
\begin{proof}
Suppose the category is connected, since the proposed type is.
Write $v$ for the large corner, of order $q\geq8$.
Deleting a retract removes at least its identity and two opposite
cross arrows, leaving at most twelve arrows and a previously
excluded category.  Assume no retract remains.

The underlying connected graph on four objects needs at least
three cross arrows.  A nontrivial strongly connected component
needs at least one further arrow.  The resulting budget
$q+3+4\leq15$ forces $q=8$ and makes all other endomorphism
monoids trivial.  Any two objects in that component then include
one with trivial endomorphism monoid; a pair of opposite arrows
exhibits it as a retract.  This is impossible.  Thus the support
is acyclic.  There are at most four cross arrows, and the other
endomorphism orders sum to at most four.

If $v$ has both an incoming and an outgoing cross arrow, those
arrows and their composite occupy a three-object chain.  Connecting
the fourth object uses the fourth cross arrow.  Therefore all
four hom-sets are singletons, $q=8$, and all other corners are
trivial.  The fourth object can attach only by an arrow from the
bottom of that chain or by an arrow to its top: attaching to its
middle, or orienting either extreme attachment the other way,
would force an additional composite hom-set.  In the first case
the bottom is initial, and in the second the top is terminal.
The category is contractible.

By opposition we may therefore suppose $v$ is a sink.  Let $D$
be the full subcategory on the other three objects.  Its corners
have orders at most two, hence rationally semisimple algebras.
If some incoming hom-set to $v$ is a singleton, then $D$ has at
most six arrows.  We claim $H^2(BD;\mathbb Q)=0$.
Unless its acyclic support contains a three-object chain, the
relative bar over the semisimple diagonal has length at most one.
If it does contain that chain, the three required cross arrows
and three identities exhaust its budget; it is the ordinary
three-element chain.  This proves the claim in both cases,
also allowing disconnected $D$.
The singleton incoming arrow is fixed by the whole large monoid,
so Lemma~\ref{C:fixed-incoming-cups} makes every degree-two
product zero.

Finally suppose every nonempty hom-set into $v$ has at least
two elements.  There can be only one neighboring object $w$:
two neighbors would use four cross arrows and still leave the
fourth object disconnected.  The remaining subcategory $D$ is
connected and needs two cross arrows.  Thus $q=8$, every corner
of $D$ is trivial, the hom-set from $w$ to $v$ has two elements,
and $D$ has exactly two singleton cross arrows.  Its support is
a three-object star.  No arrow can enter $w$, since composing it
with $w\to v$ would create another incoming hom-set to $v$.
If $w$ is the central source of the star, delete its two leaves
by the singleton cofinal deletion.  If $w$ is a source leaf of
an inward star, first delete the other source leaf and then the
central sink.  Each deletion has its singleton adjacent trivial
corner, and neither deleted vertex is adjacent to $v$.
The result is the two-object act cone with two cross arrows over
the large monoid.  These are actual homotopy-equivalence deletions.
Lemma~\ref{C:two-object-cup-pruning} gives zero products, which
again contradicts the projective-plane square.
\end{proof}
